%=====================================================================
\section{Guía de patrones --- qué tener en cuenta para el examen}
%=====================================================================
Esta sección sintetiza, a partir de los 7 finales resueltos en las secciones siguientes,
los \textbf{patrones que se repiten}: la estructura típica de cada final, los ``moldes'' de
cada ejercicio, los puntos de la consigna que suelen decidir el puntaje, y los temas de la
materia que conviene tener relacionados entre sí. No agrega contenido nuevo: es un mapa de
lo que ya está resuelto en este mismo documento.

%-----------------------------------------------------------------
\subsection{Estructura fija del final: casi siempre 3 ejercicios con roles fijos}
%-----------------------------------------------------------------
En \textbf{6 de los 7 finales} el patrón es idéntico:

\begin{center}
\small
\begin{tabular}{@{}p{0.06\textwidth}p{0.24\textwidth}p{0.28\textwidth}p{0.34\textwidth}@{}}
\toprule
\textbf{\#} & \textbf{Rol del ejercicio} & \textbf{Tema típico} & \textbf{Clases involucradas} \\
\midrule
\textbf{1} & Sistema con \textbf{cifrado simétrico} que falla en integridad/replay/modo
degradado & IoT, sensores, pulseras, dispositivos de acceso &
Primera Parte Clase 2 (Cifrado: modos ECB/CBC/CTR/stream) + Clase 3 (MAC/AEAD) + Clase 5
(nonce/replay) \\
\textbf{2} & \textbf{Protocolo criptográfico} (autenticación, reparto de secreto, firma,
SSO) con 2--3 preguntas puntuales & challenge-response, XOR ``sombras'', SSO/IdP, hash vs
HMAC & Primera Parte Clase 2/3/4/5 + Segunda Parte Clase 2 (Autenticación) \\
\textbf{3} & \textbf{Pentesting}: dar \textbf{4 hipótesis de falla} + cómo probar cada una
& app web/mobile con varias funcionalidades (pagos, IoT, exchanges, calendario, noticias) &
Segunda Parte Clase 6 (Pentesting) + a veces Clase 3 (vulnerabilidades) \\
\bottomrule
\end{tabular}
\end{center}

El único final que rompe el molde es \textbf{17/07/2023}, que reemplaza el ejercicio 2 por
una variante del mismo esquema ``sombras'' del final 2025-C2 pero con enfoque distinto ---
el patrón temático (protocolo con preguntas de robustez) se mantiene igual.

\begin{practbox}{Implicancia para estudiar}
Memorizar bien estos 3 ``moldes'' (no los enunciados exactos) cubre la gran mayoría del
examen. Practicar variantes del ejercicio 1 y tener la checklist de hipótesis de falla del
ejercicio 3 ``de memoria'' rinde más que estudiar cada final como un caso aislado.
\end{practbox}

%-----------------------------------------------------------------
\subsection{Ejercicio 1 --- el ``molde IoT/sensor con cifrado roto''}
%-----------------------------------------------------------------
Reaparece casi textual en 3 finales (pulseras Helíos / pulseras 2025-C2 / sensor--maestro
10/07/2023) con la misma trampa:
\begin{itemize}
  \item Se describe un sistema que envía $(\mathrm{ID}, \text{contador/ctr}, C)$ con
  $C=\text{cifrado simétrico}(\text{dato})$.
  \item El enunciado \textbf{nunca dice ``y esto tiene un MAC''} --- el punto es notar que
  el cifrado \emph{solo} (CTR/CBC/ECB) \textbf{da confidencialidad pero no integridad}, y
  que sin MAC un atacante puede reproducir/reenviar (\textbf{replay}) un mensaje válido
  capturado antes, o explotar la \textbf{maleabilidad} (en CTR/stream: XORear el keystream
  para voltear bits específicos sin conocer la clave; en ECB: bloques repetidos revelan
  patrones).
  \item Casi siempre hay un \textbf{``modo degradado''} (sin conexión al servidor) que baja
  la guardia de validación --- ahí la pregunta es cómo mitigar sin agregar dispositivos
  nuevos (la respuesta típica: OTP/HOTP precargadas, listas de hashes de un solo uso).
\end{itemize}

\begin{defbox}{Puntos clave a leer con cuidado en la consigna}
\begin{itemize}
  \item \textquestiondown{}El campo que viaja cifrado incluye un \textbf{contador/nonce}? Si sí, preguntar si
  el verificador chequea que sea \textbf{estrictamente creciente} (sin eso, no hay
  anti-replay real aunque el nonce exista).
  \item \textquestiondown{}Qué pasa si \textbf{falla la conexión} al servidor central? Ese es siempre el
  gancho para la pregunta de ``modo degradado''.
  \item Fijarse en el \textbf{orden de los campos dentro del cifrado} (p.\ ej.
  $\mathit{wearing}\|ctr$ vs.\ $ctr\|\mathit{wearing}$): no cambia la vulnerabilidad de
  fondo, pero el enunciado espera que la resolución sea precisa con el orden exacto.
\end{itemize}
\end{defbox}

%-----------------------------------------------------------------
\subsection{Ejercicio 2 --- el ``molde de protocolo con preguntas de robustez''}
%-----------------------------------------------------------------
Patrones que reaparecen:
\begin{itemize}
  \item \textbf{Challenge-response con hash:} $C\to S: id$, $S\to C: r$ aleatorio,
  $C\to S: H(K_c\|r)$. Casi siempre piden (a) por qué se usa $r$ y no $H(K_c)$ solo (evita
  replay), (b) si cambiar $r$ por un contador incremental es menos seguro (sí: predecible,
  aunque conceptualmente sirve como nonce único --- el problema real es previsibilidad, no
  repetición), y (c) cómo lograr autenticación \textbf{mutua sin nuevas claves} (romper la
  simetría con literales de rol, p.\ ej.\ $H(K_c\|r\|\text{``servidor''})$).
  \item \textbf{Esquema tipo ``sombras'' con XOR + hash:} reparto de un secreto en 2 partes
  (análogo a Shamir con umbral 2) donde cada parte se acompaña de un hash de la otra.
  Preguntan qué pasa si un atacante ve \textbf{solo un canal} (nada: OTP con partes
  independientes no filtra nada) y qué hace falta si ve \textbf{los dos canales} (pasar de
  hash simple a HMAC, porque el hash sin clave lo puede recalcular cualquiera que controle
  el mensaje).
  \item \textbf{SSO/IdP con firma digital:} por qué el servicio no necesita la contraseña
  (firma verificable con clave pública, no hace falta compartir secreto), cómo evitar reuso
  del token en otro servicio (agregar nonce/expiración, o firmar el campo ``servicio''
  junto con la identidad del que lo solicitó).
\end{itemize}

\begin{defbox}{Puntos clave a leer con cuidado}
\begin{itemize}
  \item Cuando el enunciado dice \textbf{``sin introducir nuevas claves''}, la respuesta
  esperada usa literales/roles dentro del hash existente, no criptografía asimétrica nueva.
  \item Distinguir siempre \textbf{integridad vs.\ confidencialidad vs.\ no repudio} --- la
  mayoría de las preguntas de ``\textquestiondown{}qué se pierde si\ldots?'' giran en torno a esta tríada
  (p.\ ej.\ MAC/HMAC da integridad+autenticidad pero \textbf{no} no-repudio; eso solo lo da
  firma digital).
\end{itemize}
\end{defbox}

%-----------------------------------------------------------------
\subsection{Ejercicio 3 --- el ``molde de pentesting: 4 hipótesis de falla''}
%-----------------------------------------------------------------
Aparece en \textbf{todos} los finales, siempre con la misma consigna implícita: dar
\textbf{4 hipótesis de falla concretas} (no genéricas) sobre un sistema descripto, y para
cada una explicar \textbf{la acción concreta} que la probaría o refutaría, siguiendo la
metodología de hipótesis de falla (Recolección de información $\to$ Hipótesis $\to$ Prueba
$\to$ Generalización $\to$ Eliminación).

\begin{warnbox}{Puntos clave a leer con cuidado en la consigna}
Esto es lo que más puntaje pierde si se ignora:
\begin{itemize}
  \item La aclaración \textbf{``no usar hipótesis genéricas''} es literal en varios
  enunciados: no alcanza con decir ``puede haber inyección SQL'' --- hay que
  \textbf{anclar la hipótesis a un campo/flujo/URL concreto} de la descripción (ej.\ ``es
  probable que el endpoint \texttt{/user\_id/public/calendar} permita enumerar
  \texttt{user\_id} y acceder a calendarios ajenos'', no ``puede haber problemas de
  autorización'').
  \item \textbf{Leer todo el enunciado buscando ``ganchos'' ya insertados por el
  enunciado}: casi siempre cada funcionalidad descripta (SMS/2FA, JWT en localStorage,
  panel de soporte, proveedor de pago externo, links de acceso temporal por mail, QR de
  vinculación, complejidad algorítmica alta) es un gancho para una hipótesis distinta. El
  enunciado suele describir 4--6 funcionalidades porque cada una es una hipótesis
  ``escondida''.
  \item Cada hipótesis necesita: (1) \textbf{dónde} está el problema (componente/flujo
  concreto), (2) \textbf{por qué es plausible} dado el diseño descripto, (3) \textbf{la
  prueba concreta} (acción ejecutable, no ``revisar el código'').
  \item Vocabulario a usar: el que está en las clases (hipótesis de falla, mediación
  completa, mínimo privilegio, relaciones de confianza, pruebas según metodología).
  \textbf{Evitar nombres de categorías estilo OWASP} (BOLA, IDOR, SSRF con esos nombres
  puntuales) si no aparecen en el material --- se puede describir el problema en esos
  términos conceptuales sin el nombre, y así queda dentro de la regla de oro.
\end{itemize}
\end{warnbox}

%-----------------------------------------------------------------
\subsection{Temas que se repiten cruzando ejercicios (para relacionar clases)}
%-----------------------------------------------------------------
\begin{itemize}
  \item \textbf{Cifrado $\neq$ integridad} (Clase 2 y 3): es el eje de casi todos los
  ejercicios 1. Todo esquema que solo cifra (ECB/CBC/CTR/stream sin MAC) es vulnerable a
  manipulación o replay --- el material lo remarca explícitamente con el ejemplo de RR.HH.
  \item \textbf{MAC da integridad + autenticidad, no no-repudio} (Clase 3): aparece cada
  vez que se compara MAC/HMAC contra firma digital.
  \item \textbf{Nonce/nonce creciente/contador vs.\ aleatorio} (Clase 5 y Segunda Parte
  Clase 2): el eje de todas las preguntas de replay y challenge-response.
  \item \textbf{Metodología de hipótesis de falla} (Segunda Parte Clase 6): un solo
  esqueleto de 5 pasos que hay que poder aplicar a \emph{cualquier} sistema descripto, sin
  conocer el sistema de antemano.
  \item \textbf{Principios de diseño (mediación completa, mínimo privilegio, defensa en
  profundidad)} (Segunda Parte Clase 3): sirven como el ``por qué'' detrás de cada
  hipótesis de falla.
\end{itemize}

%-----------------------------------------------------------------
\subsection{Huecos de material detectados al resolver}
%-----------------------------------------------------------------
\begin{ojo}
Estos términos aparecieron en enunciados de finales pero \textbf{no están, tal cual, en
los resúmenes de la materia} --- si aparecen en el examen real, resolver con los conceptos
afines que sí están (indicado entre paréntesis) y no dar por sentado vocabulario externo:
\begin{itemize}
  \item \textbf{$k$-anonimato} (usar: resistencia a preimagen, salting, efecto avalancha).
  \item \textbf{SSRF} (usar: XXE / falta de mediación completa sobre recursos internos).
  \item \textbf{ChaCha20} por nombre (usar: cifrado de flujo genérico, $E_k(m)=G(k)\oplus m$,
  PRG).
  \item \textbf{Árbol de Merkle} (usar: Merkle--Damgård es lo que sí está --- son conceptos
  distintos).
  \item \textbf{Términos OWASP puntuales} (BOLA, IDOR, ``Broken Access Control'' como
  nombre propio): usar el concepto de autorización rota / control de acceso, no el nombre
  de la categoría.
\end{itemize}
\end{ojo}

%-----------------------------------------------------------------
\subsection{Checklist rápida para el día del examen}
%-----------------------------------------------------------------
\begin{enumerate}
  \item Identificar el ``molde'' de cada ejercicio (\textquestiondown{}es IoT+cifrado? \textquestiondown{}protocolo?
  \textquestiondown{}pentesting?).
  \item Ejercicio 1: preguntarse primero ``\textquestiondown{}esto tiene MAC o no?'' y ``\textquestiondown{}qué pasa en modo
  degradado/sin conexión?''.
  \item Ejercicio 2: identificar si la pregunta es sobre integridad, confidencialidad o
  no-repudio antes de responder --- cada una tiene una herramienta distinta.
  \item Ejercicio 3: releer la descripción y anotar, funcionalidad por funcionalidad, un
  candidato a hipótesis de falla concreta (URL/campo/flujo específico) antes de escribir
  nada --- normalmente hay más ganchos en el enunciado que las 4 hipótesis pedidas, así que
  se puede elegir las 4 más sólidas.
  \item Citar explícitamente la clase/tema al fundamentar (igual que en las resoluciones de
  este documento), y si algo excede el material, decirlo en vez de inventar.
\end{enumerate}
